\section{Research Direction: Succinct Virtual Column Openings}

The Blaze compression step above uses a virtual message-oracle handle for the
folded tensor $c_\star$.  Once the verifier has sampled the coefficient vector
$\rho\in\F^{\B^\ell}$, the value at backend coordinate $b\in\B^n$ is
\[
  c_\star[b]=\sum_{h\in\B^\ell}\rho_h c[h,b].
\]
The simple implementation opens the whole committed column $c[\ast,b]$ and lets
the verifier compute this inner product.  This saves a separate commitment to
$c_\star$, but it still costs $2^\ell$ field elements for each queried backend
coordinate.

A natural research question is whether these virtual openings can be made
succinct.  Instead of sampling an arbitrary vector $\rho$, sample a row point
$r_{\mathrm{row}}\in\F^\ell$ and set
\[
  \rho_h=\chi_h(r_{\mathrm{row}})
  \qquad\text{for each }h\in\B^\ell.
\]
Then a virtual opening becomes an ordinary multilinear evaluation claim for the
opened column:
\[
  c_\star[b]
  =
  \sum_{h\in\B^\ell}\chi_h(r_{\mathrm{row}})c[h,b]
  =
  \MLE(c[\ast,b])(r_{\mathrm{row}}).
\]
So the verifier would not need to receive the whole column $c[\ast,b]$.  The
prover could instead prove the column-direction claim
\[
  \MLE(c[\ast,b])(r_{\mathrm{row}})=c_\star[b].
\]

If the backend queries coordinates $b_1,\ldots,b_q\in\B^n$, these column claims
can also be batched.  After the queried coordinates are fixed, the verifier can
sample batching weights $\eta_1,\ldots,\eta_q\in\F$ and ask for one aggregate
claim:
\[
  \MLE\left(
    h\mapsto \sum_{j=1}^q \eta_j c[h,b_j]
  \right)(r_{\mathrm{row}})
  =
  \sum_{j=1}^q \eta_j c_\star[b_j].
\]
This would replace $q$ full column openings by one proof about a
row-direction aggregate.  Combined with the usual linear-combination trick that
avoids sending all row evaluation claims $u_h$ separately, this could allow
larger values of $\ell$: the protocol would compress more rows before entering
the backend without paying $q2^\ell$ field elements in virtual openings.

The main danger is circularity.  The new column-opening proof must itself be
cheaper than opening the columns directly, and its commitment interface must be
compatible with the outer Blaze commitment.  Otherwise the protocol merely moves
the same cost into another PCS.  The soundness analysis would also need to
account for the structured choice
$\rho_h=\chi_h(r_{\mathrm{row}})$ rather than a fully arbitrary coefficient
vector.  Conceptually, though, this is the right lever: the current bottleneck is
not only the row-evaluation vector $(u_h)$, but the repeated transmission of
large committed columns used to open $c_\star$.

\subsection{Research Direction: Dimension-\(n\) Product-Tree Helpers}

The native RAA permutation check currently introduces product-tree helper
tensors
\[
  f,g\in\F^{\B^{n+1}}.
\]
The extra variable is not accidental.  The single-tensor presentation stores
both the \(2^n\) leaf factors and the \(2^n\) product-tree/internal slots in one
Boolean-addressed tensor, so the local residual has the compact form
\[
  T(1,X)-T(X,0)T(X,1).
\]
This is clean, but it forces the \(\MLEAuth\) batch that contains these helpers
to use backend dimension \(n+1\).

A possible optimization is to split the helper into two \(n\)-variate tensors:
one tensor for the leaf factors and one tensor for the internal product-tree
state.  For example, write
\[
  \ell\in\F^{\B^n},
  \qquad
  p\in\F^{\B^n},
\]
where \(\ell[b]\) stores the tagged leaf factor and \(p\) stores the product
tree layers, with one padding coordinate if needed.  The permutation check would
then need explicit relations of the form
\[
  p[\text{first internal address}]
  =
  \ell[\text{left leaf}]\,\ell[\text{right leaf}],
\]
plus the usual parent-child product relations inside \(p\), and finally the
root equality between the two product trees.

This does not reduce the total amount of helper data: two \(n\)-variate tensors
still contain \(2^{n+1}\) field elements.  The possible win is different.  If
all RAA helper objects can be expressed as \(n\)-variate tensors, then the
compiled authentication batch might avoid raising the backend dimension from
\(n\) to \(n+1\).  That would remove one backend folding layer and improve the
constant factors in the BaseFold/RMLE authentication step.

The danger is that the saved variable may reappear as relation complexity.  The
split representation would need more cases: leaf-to-first-product constraints,
internal product constraints, padding/root conventions, and terminal evaluation
requests for each relation.  The right question is therefore not whether the
same product tree can be stored in two tensors; it can.  The question is whether
the resulting relations can still be batched and authenticated without
reintroducing an \(n+1\)-variate object through the back door.
